\section{Introduction}\label{sec:introduction}
Generalized disjunctive programming (GDP) represents discrete alternatives by logical statements and associates each alternative with a set of algebraic constraints. In a convex minimization GDP, the objective and inequality functions are convex and equality constraints are affine, both globally and within each alternative; the union of alternatives is generally nonconvex. This structure suits process selection and design, because the constraints of an inactive alternative need not constrain the selected process. Logic-based outer approximation uses this structure to build mixed-integer linear masters and continuous subproblems \citep{turkay1996}. For a bounded convex disjunction, the hull formulation distributes the original variables among the alternatives and scales each copy by its selection weight \citep{balas1998,ceria1999,grossmann2003}. Convex nonlinear rows then have perspective representations, which an outer approximation can enforce through affine perspective cuts \citep{frangioni2006,bestuzheva2023}.

Even with the formulation and the family of valid inequalities fixed, the point at which a cut is constructed matters. A point tangent, as in the extended cutting plane (ECP) method \citep{westerlund1995}, linearizes at the current violated point. A radial supporting cut \citep{veinott1967}, as in the extended supporting hyperplane (ESH) method \citep{kronqvist2016}, first searches the segment from a strict interior anchor to that point and then linearizes at the boundary. The radial cut costs extra function evaluations but can remove a different part of the current relaxation. This paper asks whether that extra separation work repays its cost in a convex GDP solver.

\paragraph{Relationship to prior work.}
All ingredients of the tested solver are established, and we make no priority claim for them. \citet{serrano2020} related ESH to Kelley's method on a gauge representation, which explains why point tangents depend on the constraint representation. The SHOT solver \citep{lundell2022} already combines ESH and ECP separation with fixed-integer NLP heuristics, an ECP fallback, and single- and multiple-tree execution. For disjunctions, \citet{kronqvist2021} strengthen ESH cuts by solving separate convex problems for each term, and \citet{trespalacios2016} derive cuts from a hull relaxation strengthened by basic steps. Exact conic hull formulations \citep{bernal2024,gusev2026} and conic outer approximation \citep{coey2020} provide alternatives to finite affine outer approximation. Our study differs from these works in its question: it isolates the choice between radial and point separation under matched settings inside one solver that retains GDP term structure. \Cref{sec:related} gives the detailed comparison.

\paragraph{Approach.}
We compare ESH with ECP in one implementation, for hull and big-$M$ masters and for single-tree and repeated-master (multi-tree) execution. Within each pair, the master construction, interior initialization, incumbent recovery, validation rules and solver settings are identical. The solver uses nonlinear programming (NLP) both to find interior anchors and to recover feasible incumbents at integer selections. Its name, logic-based ESH (LB-ESH), refers to its use of the original disjunct structure, not to a new family of inequalities.

\paragraph{Contributions and main findings.}
\begin{itemize}
 \item \emph{Matched computational comparison} (\cref{sec:design,sec:results}). The study uses 51 related generated controls, 33 of them held out from the pilot split, and 27 external models. Every supplied incumbent is checked against the original model, and every failure remains in its cohort. On the held-out controls, single-tree ESH meets the numerical solve criterion on all 33 instances with either formulation; ECP does so on 32 with hull and 31 with big-$M$ (\cref{tab:primary}). Common-solved shifted geometric mean wall times favor ESH by about 4--11\% across the four primary pairs (\cref{tab:pairs}) and by 4--6\% in the single-tree pairs across three schedules. ESH generates 6--18\% fewer cuts and fewer LP iterations on average, but its mean recorded row-separation time per run is 5.5--6.6 times that of ECP (\cref{tab:work,fig:tradeoff}); node and NLP medians do not uniformly favor ESH. Formulation and tree policy have larger effects: hull/big-$M$ timing ratios are 0.468--0.652 and single/multi-tree ratios 0.676--0.893 (\cref{tab:structure}).
 \item \emph{Context from baselines and ablations} (\cref{sec:conic-results,sec:external,sec:ablations}). On the external models, ESH and ECP have equal accepted counts in every matched configuration. On the nine quadratic controls, the exact conic hull formulation is about five times faster than the prototype's hull single-tree configurations; on supported external models, the conic formulation also has lower common-solved mean times. Disabling integer-point NLP recovery reduces accepted outcomes from 69 to one of 72 pilot runs, so this implementation depends heavily on NLP recovery. Outcome-triggered baseline follow-ups are reported separately from the original runs (\cref{sec:followups}).
 \item \emph{Separation guarantees} (\cref{sec:guarantees-overview,sec:theory}). We give self-contained conditions for cut validity and separate-disjunction hulls at zero selection weights (\cref{prop:perspective,prop:poly-hull}), finite fixed-residual separation (\cref{thm:integral,thm:fractional}), approximate and finite-arithmetic cuts (\cref{prop:arithmetic}) and single-tree execution (\cref{thm:single-tree}), distinguishing valid relaxation bounds from objective-gap certificates; we claim no priority for the underlying convexity and compactness arguments. We also derive a residual-calibrated small-weight omission rule (\cref{eq:safe-skip}), which differs from the fixed cutoff used in the experiments, and a per-disjunction repair bound (\cref{prop:repair}), and show that separate-term interiors give no linear objective-error bound after intersection with global constraints (\cref{ex:intersection}).
 \item \emph{Representation and oracle-cost diagnostic} (\cref{sec:oracle}). For a fixed anchor and exact boundary roots, ESH cuts are unchanged when a row $g$ is replaced by $\phi\circ g$ with $\phi$ convex, differentiable and strictly increasing, $\phi(0)=0$ and $\phi'(0)>0$, whereas the ECP cut at the same point is never stronger (\cref{prop:invariance}), and a scalar example gives an arbitrarily long ECP transient (\cref{prop:scalar}); an explicit example shows that neither policy dominates in general (\cref{ex:nondominance}). An executable version with the frozen separator counts the function evaluations used by radial searches (\cref{sec:oracle-empirical}).
\end{itemize}

\paragraph{Scope.}
ECP shares the interior initialization even though point tangents do not need it, so the comparison measures a policy change inside this solver rather than two independently optimized implementations. The generated controls share model structures and parameters, and repeated schedules change execution order but keep the solver seed. Reported solves pass stated numerical feasibility and gap tests; they are not exact-arithmetic certificates.

\paragraph{Outline.}
\Cref{sec:related} compares this study with prior work in detail. \Cref{sec:model} defines the problem class and cuts, and \cref{sec:guarantees-overview} summarizes the guarantees proved in \cref{sec:theory}. \Cref{sec:algorithm} describes the implementation, \cref{sec:design} the experimental design, and \cref{sec:results} the results, followed by discussion and conclusions (\cref{sec:discussion,sec:conclusion}). The appendices contain the proofs, the representation diagnostic (\cref{sec:oracle}), complete model specifications (\cref{app:models}), detailed tables (\cref{app:computational}) and reproducibility information (\cref{app:reproducibility}).
